% ============================================================================
% WORKBOOK — Additional: Personal Financial Literacy
% State-specific (MN, OK, TX)
% Practice-focused reinforcement for the study guide topic
% ============================================================================
\section{Personal Financial Literacy}
\topicTitle{Personal Financial Literacy}

\begin{quickReview}{Money Management Basics}
\begin{itemize}[leftmargin=*]
    \item A \textbf{bank} holds your money in savings and checking accounts. A \textbf{credit union} is similar but member-owned (often lower fees).
    \item A \textbf{debit card} spends money you already have. A \textbf{credit card} borrows money you must pay back, usually with \textbf{interest}.
    \item \textbf{Interest rate} is the percent a lender charges (or a bank pays you) on money over time.
    \item Simple interest formula: $I = P \times r \times t$.
\end{itemize}

\medskip
\textbf{Example:} You charge $\$200$ on a credit card at $18\%$ annual interest for $1$ year. Interest $= 200 \times 0.18 \times 1 = \$36$. Total owed $= \$236$.
\end{quickReview}

% ── Warm-Up ──────────────────────────────────────────────────────────────────
\begin{practiceBox}[funGreen]{\faSun~Warm-Up}

\practiceHeader[funGreen]{Quick Interest Calculations}

\begin{multicols}{2}
\prob $\$100$ at $5\%$ for $1$ year. Find the interest. \answerBlank[3cm]
\answerExplain{$\$5$}{Use $I = Prt$: $I = 100 \times 0.05 \times 1 = \$5$. The interest earned is $\$5$.}

\prob $\$500$ at $4\%$ for $1$ year. Find the interest. \answerBlank[3cm]
\answerExplain{$\$20$}{$I = 500 \times 0.04 \times 1 = \$20$.}

\prob $\$250$ at $10\%$ for $1$ year. Find the interest. \answerBlank[3cm]
\answerExplain{$\$25$}{$I = 250 \times 0.10 \times 1 = \$25$.}

\prob $\$1{,}000$ at $3\%$ for $2$ years. Find the interest. \answerBlank[3cm]
\answerExplain{$\$60$}{$I = 1{,}000 \times 0.03 \times 2 = \$60$.}

\prob $\$800$ at $6\%$ for $1$ year. Find the total amount. \answerBlank[3cm]
\answerExplain{$\$848$}{$I = 800 \times 0.06 \times 1 = \$48$. Total $= 800 + 48 = \$848$.}
\end{multicols}
\end{practiceBox}

% ── Practice A: Comparing Payment Methods ────────────────────────────────────
\begin{practiceBox}{Debit vs.\ Credit}

\prob You buy a $\$350$ laptop with a debit card. What is the total cost? \answerBlank[3cm]
\answerExplain{$\$350$}{A debit card uses money already in your account. There is no interest, so you pay exactly $\$350$.}

\prob You buy the same $\$350$ laptop with a credit card at $18\%$ annual interest and pay it off after $1$ year. What is the total cost? \answerBlank[3cm]
\answerExplain{$\$413$}{$I = 350 \times 0.18 \times 1 = \$63$. Total $= 350 + 63 = \$413$.}

\prob How much more does the credit card cost compared to the debit card? \answerBlank[3cm]
\answerExplain{$\$63$ more}{The difference is the interest: $413 - 350 = \$63$.}

\prob You charge $\$600$ on a credit card at $20\%$ and pay it off after $6$ months. What is the interest? \answerBlank[3cm]
\answerExplain{$\$60$}{$I = 600 \times 0.20 \times 0.5 = \$60$. (Six months is half a year, so $t = 0.5$.)}

\prob Name one advantage and one disadvantage of using a credit card. \answerBlank[5cm]
\answerExplain{Advantage: buy now; Disadvantage: interest}{Credit lets you purchase before you have the cash, but you pay interest if you don't pay the full balance on time.}
\end{practiceBox}

% ── Practice B: Comparing Financial Institutions ─────────────────────────────
\begin{practiceBox}[funTeal]{Comparing Banks and Credit Unions}

\prob Bank A charges a $\$6$ monthly checking fee. Credit Union B charges $\$2$/month. How much do you save per year at the credit union? \answerBlank[3cm]
\answerExplain{$\$48$}{Bank: $6 \times 12 = \$72$. Credit union: $2 \times 12 = \$24$. Savings: $72 - 24 = \$48$.}

\prob Bank A pays $2\%$ on savings. Bank B pays $3.5\%$. You deposit $\$2{,}000$ for $1$ year. How much more does Bank B earn? \answerBlank[3cm]
\answerExplain{$\$30$ more}{Bank A: $2{,}000 \times 0.02 = \$40$. Bank B: $2{,}000 \times 0.035 = \$70$. Difference: $70 - 40 = \$30$.}

\prob A credit union pays $4\%$ on savings. You deposit $\$1{,}500$ for $2$ years. How much interest do you earn? \answerBlank[3cm]
\answerExplain{$\$120$}{$I = 1{,}500 \times 0.04 \times 2 = \$120$.}

\prob A bank charges a $\$3$ monthly savings fee but pays $2.5\%$ interest. You deposit $\$1{,}000$ for $1$ year. Do you come out ahead or behind? \answerBlank[3cm]
\answerExplain{Behind by $\$11$}{Interest earned: $1{,}000 \times 0.025 = \$25$. Fees: $3 \times 12 = \$36$. Net: $25 - 36 = -\$11$. You lose $\$11$.}

\prob A credit card charges $15\%$ annually. Another charges $22\%$. You borrow $\$400$ for $1$ year on each. How much more does the expensive card cost? \answerBlank[3cm]
\answerExplain{$\$28$ more}{Card 1: $400 \times 0.15 = \$60$. Card 2: $400 \times 0.22 = \$88$. Difference: $88 - 60 = \$28$.}
\end{practiceBox}

% ── True or False ────────────────────────────────────────────────────────────
\begin{practiceBox}[funPurple]{True or False?}

\trueOrFalse{A debit card lets you spend more than you have in your account.}
\answerExplain{False}{A debit card draws from money already in your account. If the funds are not there, the transaction is typically declined.}

\trueOrFalse{Interest on a credit card means you end up paying more than the original price.}
\answerExplain{True}{When you carry a balance on a credit card, the lender charges interest, so the total cost is higher than the purchase price.}

\trueOrFalse{A credit union is owned by its members and often has lower fees than a bank.}
\answerExplain{True}{Credit unions are member-owned cooperatives that tend to offer lower fees and better savings rates than traditional banks.}

\trueOrFalse{Paying your credit card bill in full every month avoids interest charges.}
\answerExplain{True}{Most credit cards do not charge interest if the full balance is paid by the due date each month.}
\end{practiceBox}

% ── Word Problems ────────────────────────────────────────────────────────────
\begin{practiceBox}[funBlue]{\faGlobe~Word Problems}

\wordProblem{Tanya opens a savings account with $\$750$ at $3\%$ simple interest. How much will she have after $3$ years?}{}
\answerExplain{$\$817.50$}{$I = 750 \times 0.03 \times 3 = \$67.50$. Total $= 750 + 67.50 = \$817.50$.}

\wordProblem{Marcus uses a credit card to buy a $\$500$ phone and pays it back after $2$ years at $16\%$ annual interest. What is the total amount he pays?}{}
\answerExplain{$\$660$}{$I = 500 \times 0.16 \times 2 = \$160$. Total $= 500 + 160 = \$660$.}
\end{practiceBox}

% ── Challenge ────────────────────────────────────────────────────────────────
\begin{challengeBox}
\prob Credit Card A charges $18\%$ interest with no annual fee. Credit Card B charges $12\%$ interest but has a $\$50$ annual fee. If you carry a $\$1{,}000$ balance for $1$ year, which card costs less total? \answerBlank[3cm]
\answerExplain{Card B costs less ($\$170$ vs.\ $\$180$)}{Card A interest: $1{,}000 \times 0.18 = \$180$. Card B interest plus fee: $1{,}000 \times 0.12 + 50 = 120 + 50 = \$170$. Card B is $\$10$ cheaper.}
\end{challengeBox}

\encouragement{Awesome job learning about money management! Smart financial choices start with understanding how interest works!}
